% Archivo: figuras/graficos-likert.tex
% Gráficos de resultados de evaluación Likert

% Cargar datos
% Archivo: datos/likert-data.tex
% Datos completos de evaluación SocialScrum

% Datos generales (todas las preguntas)
\pgfplotstableread{
   Pregunta E5 E4 E3 E2 E1
   P1 8 5 2 0 0
   P2 6 6 2 1 0
   P3 4 5 4 2 0
   P4 7 6 2 0 0
   P5 9 4 2 0 0
   P6 5 4 4 2 0
   P7 10 4 1 0 0
   P8 9 5 1 0 0
   P9 8 5 2 0 0
   P10 6 7 2 0 0
   P11 4 6 3 2 0
   P12 5 5 3 2 0
   P13 8 6 1 0 0
   P14 6 5 3 1 0
   P15 7 6 2 0 0
}\datatablalikert

% Datos por criterio: CLARIDAD (P1, P2, P3)
\pgfplotstableread{
   Pregunta E5 E4 E3 E2 E1
   P1 8 5 2 0 0
   P2 6 6 2 1 0
   P3 4 5 4 2 0
}\dataclaridad

% Datos por criterio: COMPLETITUD (P4, P5, P6)
\pgfplotstableread{
   Pregunta E5 E4 E3 E2 E1
   P4 7 6 2 0 0
   P5 9 4 2 0 0
   P6 5 4 4 2 0
}\datacompletitud

% Datos por criterio: IDONEIDAD (P7, P8, P9)
\pgfplotstableread{
   Pregunta E5 E4 E3 E2 E1
   P7 10 4 1 0 0
   P8 9 5 1 0 0
   P9 8 5 2 0 0
}\dataidoneidad

% Datos por criterio: FACILIDAD DE USO (P10, P11, P12)
\pgfplotstableread{
   Pregunta E5 E4 E3 E2 E1
   P10 6 7 2 0 0
   P11 4 6 3 2 0
   P12 5 5 3 2 0
}\datafacilidaduso

% Datos por criterio: AGILIDAD (P13, P14, P15)
\pgfplotstableread{
   Pregunta E5 E4 E3 E2 E1
   P13 8 6 1 0 0
   P14 6 5 3 1 0
   P15 7 6 2 0 0
}\dataagilidad

% Datos generales pruebas (todas las preguntas)
\pgfplotstableread{
   Pregunta E5 E4 E3 E2 E1
   P1 3 3 1 1 0
   P2 3 2 2 1 0
   P3 4 4 0 0 0
   P4 7 1 0 0 0
   P5 8 0 0 0 0
   P6 5 2 1 0 0
   P7 3 4 1 0 0
   P8 2 4 2 0 0
   P9 8 0 0 0 0
   P10 6 1 1 0 0
   P11 1 6 1 0 0
   P12 5 0 3 0 0
   P13 4 2 1 1 0
   P14 3 2 3 0 0
   P15 2 2 2 2 0
}\datatablalikertpruebas

% ============================================================
% FIGURA 1: Distribución general (todas las preguntas)
% ============================================================
\begin{figure}[H]
    \centering
    \begin{tikzpicture}
        \begin{axis}[
                xbar stacked,
                width=0.95\textwidth,
                height=14cm,
                xlabel={Número de respuestas},
                ylabel={Preguntas del instrumento},
                symbolic y coords={P15,P14,P13,P12,P11,P10,P9,P8,P7,P6,P5,P4,P3,P2,P1},
                ytick=data,
                legend style={at={(0.5,-0.08)}, anchor=north, legend columns=5, font=\footnotesize},
                bar width=9pt,
                xmin=0, xmax=15,
                enlarge y limits=0.03,
            ]
            \addplot[fill=blue!80, draw=black, line width=0.3pt] table[x=E5, y=Pregunta] {\datatablalikert};
            \addlegendentry{Muy satisfecho (5)}
            \addplot[fill=blue!50, draw=black, line width=0.3pt] table[x=E4, y=Pregunta] {\datatablalikert};
            \addlegendentry{Satisfecho (4)}
            \addplot[fill=gray!50, draw=black, line width=0.3pt] table[x=E3, y=Pregunta] {\datatablalikert};
            \addlegendentry{Neutral (3)}
            \addplot[fill=orange!70, draw=black, line width=0.3pt] table[x=E2, y=Pregunta] {\datatablalikert};
            \addlegendentry{Poco satisfecho (2)}
            \addplot[fill=red!70, draw=black, line width=0.3pt] table[x=E1, y=Pregunta] {\datatablalikert};
            \addlegendentry{Nada satisfecho (1)}
        \end{axis}
    \end{tikzpicture}
    \caption{Distribución general de respuestas del instrumento de evaluación según escala Likert}
    \label{fig:likert-general}
\end{figure}

% ============================================================
% FIGURA 2: Criterio CLARIDAD
% ============================================================
\begin{figure}[H]
    \centering
    \begin{tikzpicture}
        \begin{axis}[
                xbar stacked,
                width=0.85\textwidth,
                height=6cm,
                xlabel={Número de respuestas},
                ylabel={Preguntas},
                symbolic y coords={P3,P2,P1},
                ytick=data,
                legend style={at={(0.5,-0.15)}, anchor=north, legend columns=5, font=\footnotesize},
                bar width=12pt,
                xmin=0, xmax=15,
                enlarge y limits=0.1,
            ]
            \addplot[fill=blue!80, draw=black, line width=0.3pt] table[x=E5, y=Pregunta] {\dataclaridad};
            \addlegendentry{Muy satisfecho (5)}
            \addplot[fill=blue!50, draw=black, line width=0.3pt] table[x=E4, y=Pregunta] {\dataclaridad};
            \addlegendentry{Satisfecho (4)}
            \addplot[fill=gray!50, draw=black, line width=0.3pt] table[x=E3, y=Pregunta] {\dataclaridad};
            \addlegendentry{Neutral (3)}
            \addplot[fill=orange!70, draw=black, line width=0.3pt] table[x=E2, y=Pregunta] {\dataclaridad};
            \addlegendentry{Poco satisfecho (2)}
            \addplot[fill=red!70, draw=black, line width=0.3pt] table[x=E1, y=Pregunta] {\dataclaridad};
            \addlegendentry{Nada satisfecho (1)}
        \end{axis}
    \end{tikzpicture}
    \caption{Distribución de respuestas para el criterio de Claridad}
    \label{fig:likert-claridad}
\end{figure}

% ============================================================
% FIGURA 3: Criterio COMPLETITUD
% ============================================================
\begin{figure}[H]
    \centering
    \begin{tikzpicture}
        \begin{axis}[
                xbar stacked,
                width=0.85\textwidth,
                height=6cm,
                xlabel={Número de respuestas},
                ylabel={Preguntas},
                symbolic y coords={P6,P5,P4},
                ytick=data,
                legend style={at={(0.5,-0.15)}, anchor=north, legend columns=5, font=\footnotesize},
                bar width=12pt,
                xmin=0, xmax=15,
                enlarge y limits=0.1,
            ]
            \addplot[fill=blue!80, draw=black, line width=0.3pt] table[x=E5, y=Pregunta] {\datacompletitud};
            \addlegendentry{Muy satisfecho (5)}
            \addplot[fill=blue!50, draw=black, line width=0.3pt] table[x=E4, y=Pregunta] {\datacompletitud};
            \addlegendentry{Satisfecho (4)}
            \addplot[fill=gray!50, draw=black, line width=0.3pt] table[x=E3, y=Pregunta] {\datacompletitud};
            \addlegendentry{Neutral (3)}
            \addplot[fill=orange!70, draw=black, line width=0.3pt] table[x=E2, y=Pregunta] {\datacompletitud};
            \addlegendentry{Poco satisfecho (2)}
            \addplot[fill=red!70, draw=black, line width=0.3pt] table[x=E1, y=Pregunta] {\datacompletitud};
            \addlegendentry{Nada satisfecho (1)}
        \end{axis}
    \end{tikzpicture}
    \caption{Distribución de respuestas para el criterio de Completitud}
    \label{fig:likert-completitud}
\end{figure}

% ============================================================
% FIGURA 4: Criterio IDONEIDAD
% ============================================================
\begin{figure}[H]
    \centering
    \begin{tikzpicture}
        \begin{axis}[
                xbar stacked,
                width=0.85\textwidth,
                height=6cm,
                xlabel={Número de respuestas},
                ylabel={Preguntas},
                symbolic y coords={P9,P8,P7},
                ytick=data,
                legend style={at={(0.5,-0.15)}, anchor=north, legend columns=5, font=\footnotesize},
                bar width=12pt,
                xmin=0, xmax=15,
                enlarge y limits=0.1,
            ]
            \addplot[fill=blue!80, draw=black, line width=0.3pt] table[x=E5, y=Pregunta] {\dataidoneidad};
            \addlegendentry{Muy satisfecho (5)}
            \addplot[fill=blue!50, draw=black, line width=0.3pt] table[x=E4, y=Pregunta] {\dataidoneidad};
            \addlegendentry{Satisfecho (4)}
            \addplot[fill=gray!50, draw=black, line width=0.3pt] table[x=E3, y=Pregunta] {\dataidoneidad};
            \addlegendentry{Neutral (3)}
            \addplot[fill=orange!70, draw=black, line width=0.3pt] table[x=E2, y=Pregunta] {\dataidoneidad};
            \addlegendentry{Poco satisfecho (2)}
            \addplot[fill=red!70, draw=black, line width=0.3pt] table[x=E1, y=Pregunta] {\dataidoneidad};
            \addlegendentry{Nada satisfecho (1)}
        \end{axis}
    \end{tikzpicture}
    \caption{Distribución de respuestas para el criterio de Idoneidad}
    \label{fig:likert-idoneidad}
\end{figure}

% ============================================================
% FIGURA 5: Criterio FACILIDAD DE USO
% ============================================================
\begin{figure}[H]
    \centering
    \begin{tikzpicture}
        \begin{axis}[
                xbar stacked,
                width=0.85\textwidth,
                height=6cm,
                xlabel={Número de respuestas},
                ylabel={Preguntas},
                symbolic y coords={P12,P11,P10},
                ytick=data,
                legend style={at={(0.5,-0.15)}, anchor=north, legend columns=5, font=\footnotesize},
                bar width=12pt,
                xmin=0, xmax=15,
                enlarge y limits=0.1,
            ]
            \addplot[fill=blue!80, draw=black, line width=0.3pt] table[x=E5, y=Pregunta] {\datafacilidaduso};
            \addlegendentry{Muy satisfecho (5)}
            \addplot[fill=blue!50, draw=black, line width=0.3pt] table[x=E4, y=Pregunta] {\datafacilidaduso};
            \addlegendentry{Satisfecho (4)}
            \addplot[fill=gray!50, draw=black, line width=0.3pt] table[x=E3, y=Pregunta] {\datafacilidaduso};
            \addlegendentry{Neutral (3)}
            \addplot[fill=orange!70, draw=black, line width=0.3pt] table[x=E2, y=Pregunta] {\datafacilidaduso};
            \addlegendentry{Poco satisfecho (2)}
            \addplot[fill=red!70, draw=black, line width=0.3pt] table[x=E1, y=Pregunta] {\datafacilidaduso};
            \addlegendentry{Nada satisfecho (1)}
        \end{axis}
    \end{tikzpicture}
    \caption{Distribución de respuestas para el criterio de Facilidad de Uso}
    \label{fig:likert-facilidad}
\end{figure}

% ============================================================
% FIGURA 6: Criterio AGILIDAD
% ============================================================
\begin{figure}[H]
    \centering
    \begin{tikzpicture}
        \begin{axis}[
                xbar stacked,
                width=0.85\textwidth,
                height=6cm,
                xlabel={Número de respuestas},
                ylabel={Preguntas},
                symbolic y coords={P15,P14,P13},
                ytick=data,
                legend style={at={(0.5,-0.15)}, anchor=north, legend columns=5, font=\footnotesize},
                bar width=12pt,
                xmin=0, xmax=15,
                enlarge y limits=0.1,
            ]
            \addplot[fill=blue!80, draw=black, line width=0.3pt] table[x=E5, y=Pregunta] {\dataagilidad};
            \addlegendentry{Muy satisfecho (5)}
            \addplot[fill=blue!50, draw=black, line width=0.3pt] table[x=E4, y=Pregunta] {\dataagilidad};
            \addlegendentry{Satisfecho (4)}
            \addplot[fill=gray!50, draw=black, line width=0.3pt] table[x=E3, y=Pregunta] {\dataagilidad};
            \addlegendentry{Neutral (3)}
            \addplot[fill=orange!70, draw=black, line width=0.3pt] table[x=E2, y=Pregunta] {\dataagilidad};
            \addlegendentry{Poco satisfecho (2)}
            \addplot[fill=red!70, draw=black, line width=0.3pt] table[x=E1, y=Pregunta] {\dataagilidad};
            \addlegendentry{Nada satisfecho (1)}
        \end{axis}
    \end{tikzpicture}
    \caption{Distribución de respuestas para el criterio de Agilidad}
    \label{fig:likert-agilidad}
\end{figure}